% !TeX root = ../document_maitre.tex

\textbf{Lien :} \url{https://github.com/floflo-boop/Livrable-technique}\newline

Ce dépôt \textit{Github} se constitue de deux documents PDF. Ce sont des documents produits dans le cadre du stage à propos de la création de la fonctionnalité de \guillemet{recherche augmentée}.\footnote{Ces documents professionnels sont la propriété de l'entreprise et ne peuvent être diffusés librement} L'objectif de ce dépôt est de montrer ce qui a été mis en place pour la fonctionnalité de \guillemet{recherche avancée} bien que cela ne soit pas exhaustif.\newline

Le premier document \textit{spécification.pdf} contient la spécification rédigée afin d'établir la fonctionnalité de la \guillemet{recherche augmentée}. Elle explicite les attendus et permet d'établir les règles et fonctions à mettre en place pour l'expérience et l'interface utilisateur.

Loin d'être un document \guillemet{technique} au sens qu'il ne développe pas de processus, ni de code, son objectif est de contraindre le développement de la fonctionnalité par un raisonnement centré sur l'adoption de l'outil par l'utilisateur. C'est donc un document de contrainte afin de réaliser les attentes utilisateurs dans une fonctionnalité, et non l'inverse.\newline

Le second document \textit{index.pdf} est un index \gls{ES} commenté pour être utilisé par la fonctionnalité de \guillemet{recherche augmentée}. L'objectif est de montrer comment ce structure un index \gls{ES} et comment les exploiter pour cette fonctionnalité. Il faut garder en tête que chaque index doit être commenté de cette manière afin de faire fonctionner la \guillemet{recherche augmentée}.